\scanpage{92}
\mspage{- 174 -}
\hfill IX

(13.3.9) correspond simplement à l'homomorphisme
\[
E^t/R_{\unsure{o|0}S}^t=E^t/R^t\longrightarrow E/R_{oS}=E/R^f
\]
induit par les inclusions
\[
E^t\hookrightarrow E,\quad R^t\hookrightarrow R^f\quad .
\]

Utilisant le fait que (13.3.9) induit une immersion sur les fibres en $s$, on en conclut que l'homomorphisme
\[
R_o/R_o^t\longrightarrow E_o/E_o^t=(E/E^t)_o
\]
est injectif, d'où il résulte aussitôt que
\[
R_o/R_o^t\longrightarrow(E/E^t)_{s'}
\]
est encore injectif pour $s'$ voisin de $s$, ce qui implique que, quitte à rapetisser $S$ autour de $s$, on peut supposer que (13.3.9) \uline{est une immersion fermée}.

Il s'impose alors d'introduire
\[
B=A^{\handwritten{\natural}}/T\simeq(E/E^t)/(R^f/R^t)\quad , \tag*{(13.3.12)}
\]
de sorte qu'on obtient sur $A^{\handwritten{\natural}}$ une structure d'extension
\[
0\longrightarrow T\longrightarrow A^{\handwritten{\natural}}\longrightarrow B\longrightarrow 0\quad , \tag*{(13.3.13)}
\]
$B$ étant un groupe analytique relatif sur $S$, lisse et séparé sur $S$ et à fibres connexes. Par construction même, on a donc
\[
\underline{\mathrm{Lie}}(B)\simeq\underline E/\underline E^t\quad , \tag*{(13.3.14)}
\]
et la suite exacte exponentielle pour $B$

% Faint, partly covered numerals at lower strip. Attribution uncertain.
\margin{\unsure{492|1192}}
% printed label at foot
492
% archive sheet number: 91
